\chapter{Censoring}
\label{ch:censoring}

\section{Censoring Mode}

Activate censoring with the \texttt{censoring=true} class option:

\begin{minted}{latex}
\documentclass[censoring=true]{omnilatex}
\end{minted}

This loads the \texttt{censor} package and enables all censoring commands
described below.  When \texttt{censoring=false} (the default), all censoring
commands are no-ops and produce no output.

\section{\textbackslash censor\{text\}}

The \cs{censor} command applies inline black-bar redaction:

\begin{minted}{latex}
The secret project codename is \censor{Operation Sunflower}.
\end{minted}

The redacted text is replaced with a black rectangle matching the text width.
Use \cs{StopCensoring} to disable censoring from a point onward:

\begin{minted}{latex}
\censor{Classified information}
\StopCensoring
This text is no longer censored.
\end{minted}

\section{\textbackslash todo\{text\}}

The \cs{todo} command inserts a marginal note for internal review:

\begin{minted}{latex}
\todo{Verify these figures with the experimental team.}
\end{minted}

TODO notes are rendered only when the \texttt{todonotes} class option is
enabled.  In production builds they are automatically suppressed, ensuring no
draft annotations leak into the final PDF.

\section{\textbackslash censorbox\{\ldots\}}

For larger redactions that span multiple lines, use \cs{censorbox}:

\begin{minted}{latex}
\censorbox{%
  The following dataset contains personally identifiable
  information (PII) and must not be published:
  \begin{itemize}
    \item Names, addresses, phone numbers
    \item Medical records
    \item Financial account details
  \end{itemize}
}
\end{minted}

The box renders as a solid black rectangle covering the entire content area.

\section{Practical Use}

Censoring is designed for two primary workflows:

\begin{description}
  \item[Thesis Submission] Advisors and reviewers often insert comments,
        suggestions, and revisions during the drafting process.  Before final
        submission, enable \texttt{censoring=true} to redact all reviewer
        annotations while preserving the document structure.

  \item[Sensitive Data Redaction] Research papers handling PII, proprietary
        data, or classified material can use censoring commands to prepare
        public-facing versions without maintaining separate source files.
\end{description}

Both workflows benefit from the fact that censoring is a compile-time
feature---no manual editing of source content is required.
